\documentclass{article}
\usepackage[margin=2cm]{geometry}
\usepackage{amsmath}
\usepackage{graphicx}
\usepackage{booktabs}
\usepackage[utf8]{inputenc}
\usepackage{ragged2e}
\setlength{\emergencystretch}{3em}
\begin{document}
\sloppy

\noindent\hfill\includegraphics[height=1.8em]{assets/figure_header_logo.png}\par\noindent\hrule\vspace{4pt}

The non-zero components of the Ricci tensor are


\begin{equation}
\begin{array}{c c c}{{R_{\mathrm{out}}=-5\left(f^{\prime\prime}+(f^{\prime})^{2}\right)}}&{{\qquad\qquad}}&{{R_{\mathrm{mn}}=-g_{\mathrm{mn}}\left(f^{\prime\prime}+5(f^{\prime})^{2}\right)}}\\ {{\cdot}}&{{\cdot}}&{{\cdot}}&{{\cdot}}&{{\qquad,}}\end{array}\qquad\qquad\qquad}}&{{\qquad}\end{array
\end{equation}


while the Ricci scalar is given by


\begin{equation}
R=-10f^{\prime\prime}-30(f^{\prime})^{2}~ (2)
\end{equation}


Furthermore, we have that G = e5f g, where g is the determinant of the unit S metric.
Upon integration by parts, part of the Einstein-Hilbert term cancels with the Gibbons-
Hawking term to give the following simple expression


\begin{equation}
\begin{array}{l l l l l l l}{{S=\int d u\int d^{5}x\sqrt{g}\,e^{5f}\left[-5\left(\left(f^{\prime}\right)^{2}+4W^{2}\right)+2\bar{Z}_{\mathrm{hin}}\right]}}&{{}}&{{}}&{{(3)}}\\ {{..}}&{{}}&{{.}}&{{}}&{{.}}&{{}}&{{}}&{{.}}&{{}}\end{array}
\end{equation}


The restriction to the flat case was not strictly necessary so far, but it will be crucial in the
next step. The gradient flow equations, together with the chain-rule, allows us to rewrite


\begin{equation}
\mathcal{L}_{\mathrm{kin}}=-2\,W^{\prime}\qquad (4)
\end{equation}


Plugging this into and using the BPS equation of the warp factor, we find


\begin{equation}
S=-4\int d^{5}x\sqrt{g}\,e^{5f}W\vert_{0}^{\Lambda}\qquad
\end{equation}


where A is the UV cutoff. Only the A part of the action contributes, since e5fW|o vanishes.
due to the close-off of the geometry. Removing the UV cutoff A  oo is equivalent to
removing the cutoff e on our asymptotic coordinate z, i.e. e -> 0. From the UV asymptotics
we find that in this limit the factor e5f diverges like


\begin{equation}
\begin{array}{c c}{{\displaystyle{\epsilon^{5f}\sim{\frac{1}{\varepsilon^{5}}}\sim{\frac{1}{\varepsilon^{5}}}\qquad (6)\qquad}}\end{array}}
\end{equation}


tial are relevant for obtaining counterterms. All the higher-order terms vanish as the cutoff
is removed. We may thus legitimately insert the approximate superpotential into to get the
counterterms,


\begin{equation}
S_{\mathrm{ct}}^{(W)}=4\int d^{6}x\sqrt{\gamma}\left[\frac{1}{2}+\frac{3}{4}\sigma^{2}+\frac{1}{16}(\phi^{0})^{2}-\frac{3}{16}(\phi^{0})^{2}+\frac{1}{192}(\phi^{0})^{4}-\frac{3}{16}(\phi^{0})^{2}\sigma\right]\quad\quad\quad(7)
\end{equation}


where  is the induced metric on the z = e boundary. All fields are evaluated at z = e. This
gives all finite and infinite counterterms for the flat domain wall solutions..

\end{document}
